% Independent derivation, 2026-09-23. SRC1--SRC15.
\section{Geometric realization of the full supported restricted adelic object}
\label{sec:supported-realization-comparison}

\paragraph{Sources and scope.}
The original restricted simplicial object is due to Alain Connes and
Caterina Consani, \href{https://arxiv.org/abs/2205.01391v2}{\emph{Riemann--Roch
for $\overline{\operatorname{Spec}\mathbb Z}$}}, original author source
\nolinkurl{RR-J-final.tex}, Appendix B, particularly
\texttt{descript1}, \texttt{descript2}, \texttt{descript1sub},
\texttt{descript1sub1}, \texttt{Maydefn}, and \texttt{proppi0}.
Their divisor restrictions are in \texttt{setup}--\texttt{setup1}.
The support carrier is that of The Clankers,
\href{https://github.com/KokunoYumeto/zeta-function-research-reader/blob/a2851f9eb352e7e4376bc629de86fa4ed9c1c434/workbenches/splitzero-tandem/foundations/split-support-geometry-v11/split_support_geometry_arithmetic_curve_v11.tex#L1226}
{\emph{Split Support Geometry}, version 11, Definition \texttt{def:lattice-split}}.
The comparison proved here applies to the full independent-leg extension
in ASC1--ASC5b. It concerns geometric realizations of simplicial sets;
the adele amplitudes have the discrete topology at this stage, as in
that simplicial-set construction. No topology on $L$ is assumed.
In particular it does not identify realization homotopy groups with
the strict spherical-simplex relation used in the original source.

\subsection{The two original arrays and their amplitude map}

Let $L$ be a bounded distributive lattice with $0_L\ne1_L$, and set
\begin{equation}\tag{SRC1}
\begin{gathered}
G_L(A)=\{(0,\lambda):\lambda\in L\}\cup(A\times\{1_L\}),\qquad
(x,\lambda)+(y,\mu)=(x+y,\lambda\vee\mu),\\
\tau_A=(0,0_L),\qquad e_A=(0,1_L),\\
D=\sum_p d_p[p]+d_\infty[\infty],\qquad
N=\prod_p p^{-d_p},\qquad a=e^{d_\infty},\qquad
O_f(D)=N\widehat{\mathbb Z}.
\end{gathered}
\end{equation}
Here $N,a>0$ and only finitely many $d_p$ are nonzero. Put
$B_L=G_L(\mathbb A_{\mathbb Q})$ and
$P_L=G_L(\mathbb Q)\times G_L(\mathbb A_{\mathbb Q})$; its two
factors have independent labels. The action map is
\begin{equation}\tag{SRC2}
\Phi_L((q,\alpha),(o,\beta))=(q+o,\alpha\vee\beta),
\end{equation}
where $q$ is embedded diagonally in the adeles.
For a pointed finite set $r_+$ the supported $n$-simplices are
\begin{equation}\tag{SRC3}
\begin{gathered}
b=(b_j)\in B_L^r,\qquad
z_\ell=(z_{\ell j})\in P_L^r\quad(1\le\ell\le n),\\
z_{\ell j}=((q_{\ell j},\alpha_{\ell j}),
             (o_{\ell j},\beta_{\ell j})),\qquad
o_{\ell j,f}\in O_f(D),\qquad
\sum_{\ell=1}^n\sum_{j=1}^r|o_{\ell j,\infty}|\le a.
\end{gathered}
\end{equation}
The amplitude object $Y_r$ has exactly these arrays, finite
restrictions and total bound, with the labels forgotten. Write $X_r$
for the supported object. Neither object replaces the total bound
by individual bounds on the entries.

For the proof temporarily use the usual chronological bar faces
$\partial_i$:
\begin{equation}\tag{SRC4}
\begin{aligned}
\partial_0(b;z_1,\ldots,z_n)
  &=(b+\Phi_Lz_1;z_2,\ldots,z_n),\\
\partial_i(b;z_1,\ldots,z_n)
  &=(b;z_1,\ldots,z_i+z_{i+1},\ldots,z_n)
      &&(0<i<n),\\
\partial_n(b;z_1,\ldots,z_n)&=(b;z_1,\ldots,z_{n-1}).
\end{aligned}
\end{equation}
The degeneracy $\sigma_i$ inserts $(\tau_{\mathbb Q},\tau_{\mathbb A})$
after the first $i$ arrows. ASC5 uses the reversed convention
$d_i=\partial_{n-i}$ and $s_i=\sigma_{n-i}$. Reversal of barycentric
coordinates identifies the two realizations, including their pointed
finite-set maps, so this convention retains the actual ASC object.
All faces preserve (SRC3) by the triangle inequality, and degeneracies
insert real component zero. Summation over the nonbasepoint fibres of
a pointed finite-set map preserves the same bound and joins each
leg's labels separately. An empty fibre inserts the relevant $\tau$.

Thus forgetting labels gives the original pointed simplicial
$\Gamma$-map
\begin{equation}\tag{SRC5}
p_D:X(D)\longrightarrow Y(D).
\end{equation}
It is onto in every degree, but this is not the argument for the
homotopy assertion below. Over the all-zero array its actual
degree-$(n,r)$ fibre is $L^{(2n+1)r}$; the maps in that fibre are
the prescribed joins and identity insertions.

\subsection{The zero-cost arrow and its exact bar prism}

Fix $r$. Let $c=c_r\in P_L^r$ have each coordinate
$(e_{\mathbb Q},e_{\mathbb A})$. Then
\begin{equation}\tag{SRC6}
c+c=c,\qquad \Phi_Lc=e_{\mathbb A}^r,\qquad
c_{\ell j,\mathrm{amplitude}}=0.
\end{equation}
It is admitted with real cost zero for every original divisor $D$.
For $r=0$ this is the unique empty array.
Define a map $t_r:Y_r\to X_r$ on simplices by attaching top labels
to the vertex and to each of the two legs of every arrow, including
amplitude-zero entries. Then
\begin{equation}\tag{SRC7}
p_rt_r=\operatorname{id},\qquad
t_rp_r(b;z_1,\ldots,z_n)
 =(b+\Phi_Lc;z_1+c,\ldots,z_n+c).
\end{equation}
Both maps in (SRC7) commute with every face. For the initial face,
the supported vertex after the action is
$b+\Phi_Lc+\Phi_L(z_1+c)=b+\Phi_Lz_1+\Phi_Lc$.
For an interior face, $(z_i+c)+(z_{i+1}+c)=z_i+z_{i+1}+c$.
The last face only drops an arrow. These identities also cover
zero amplitudes and independent zero-leg labels.

The map $t_r$ is a map of semisimplicial sets, not generally of
simplicial sets: $t_r$ turns an amplitude identity arrow into $c$,
whereas a degeneracy in $X_r$ inserts
$(\tau_{\mathbb Q},\tau_{\mathbb A})$. For $r>0$, the two arrows
are different. No simplicial section has been assumed.

For $0\le i\le n$ define the following admitted $(n+1)$-simplex:
\begin{equation}\tag{SRC8}
H_i(b;z_1,\ldots,z_n)
 =(b;z_1,\ldots,z_i,c,z_{i+1}+c,\ldots,z_n+c).
\end{equation}
Every rational and adelic amplitude in this array is the original
amplitude, with one new zero row. The finite components still belong
to $N\widehat{\mathbb Z}$ and the full real sum in (SRC3) is unchanged,
not merely bounded above by a different estimate. The two source
labels of every displayed $z_\ell+c$ are joined independently with
$1_L$. They have not first been identified with one common label.

The prism identities are
\begin{equation}\tag{SRC9}
\begin{aligned}
\partial_0H_0&=t_rp_r,&
\partial_{n+1}H_n&=\operatorname{id},\\
\partial_jH_i&=H_{i-1}\partial_j &&(j<i),\\
\partial_jH_i&=H_i\partial_{j-1} &&(j>i+1),\\
\partial_iH_i&=\partial_iH_{i-1} &&(1\le i\le n).
\end{aligned}
\end{equation}
Here all operators on the right have the degree determined by their
input. For $j<i$ the deleted or merged arrows precede the inserted
$c$; the case $j=0$ updates the same vertex by the same $\Phi_Lz_1$.
For $j>i+1$ they follow $c$; merging uses
$(z+c)+(z'+c)=z+z'+c$. At the seam the two results are
$z_i+c$ and $c+(z_i+c)$, which agree. The endpoint identities follow
by dropping the last $c$, or by applying $\Phi_Lc$ to the initial
vertex. This proves every identity without an additive inverse in
the supported monoids.

Let $\|K\|$ denote fat realization, in which only face identifications
are imposed. Triangulate $\Delta^n\times[0,1]$ by the $n+1$ simplices
with vertices
\begin{equation}\tag{SRC10}
(v_0,0),\ldots,(v_i,0),(v_i,1),\ldots,(v_n,1)
       \qquad(0\le i\le n).
\end{equation}
On that simplex use the characteristic map labelled by $H_i$.
The last three identities in (SRC9) make these maps agree on all
common faces and with all face identifications from the original
$n$-simplex. The endpoint identities give a continuous homotopy
from the identity to $\|t_r\|\|p_r\|$. Continuity follows on each
prism and hence on the quotient: product with the locally compact
Hausdorff interval preserves the quotient map defining realization.
Since $p_rt_r$ is the identity,
\begin{equation}\tag{SRC11}
\|p_r\|:\|X_r\|\longrightarrow\|Y_r\|
\quad\hbox{is a homotopy equivalence.}
\end{equation}
This argument does not pretend that (SRC8) respects degeneracies.
After applying $p_r$, however, its inserted arrow is the amplitude
identity. Consequently its image in the ordinary realization
$|Y_r|$ is constant in the homotopy parameter. The displayed
homotopy lies over the fixed map $\|X_r\|\to|Y_r|$.

\subsection{A cellular proof of the fat-to-ordinary comparison}

For completeness the following argument proves the realization fact
needed here, so no unverified theorem about a simplicial topological
space is imported. For any simplicial set $K$, there is a natural
quotient
\begin{equation}\tag{SRC12}
q_K:\|K\|\longrightarrow|K|
\end{equation}
which also imposes the degeneracy identifications. It is a weak
homotopy equivalence by the complete cellular argument below.
The further conclusion that it is a homotopy equivalence uses the
classical CW Whitehead theorem, originating in J.\ H.\ C.\ Whitehead,
\href{https://doi.org/10.1090/S0002-9904-1949-09175-9}
{\emph{Combinatorial homotopy. I}},
\emph{Bulletin of the American Mathematical Society} \textbf{55}
(1949), 213--245. This is an attribution of the classical theorem;
the present source-reading claim does not include the 1949 article.

First, every simplex has a unique expression $\alpha^*x$ where
$x$ is nondegenerate and $\alpha$ is an order-preserving surjection.
Existence follows by repeatedly removing a degeneracy. To verify
uniqueness, the indices $i$ for which a simplex lies in the image
of $\sigma_i$ are exactly the indices with
$\alpha(i)=\alpha(i+1)$. One inclusion follows by factoring
$\alpha$ through the elementary surjection identifying $i,i+1$.
For the converse, if $\alpha(i)<\alpha(i+1)$, these are consecutive
values, and one can choose an order-preserving section of $\alpha$
containing both $i$ and $i+1$. Pulling an expression in the image
of $\sigma_i$ back along that section would express $x$ as a
degeneracy, a contradiction. These equal-adjacent indices determine
$\alpha$ uniquely. Pullback along any section then determines $x$.

Fat realization is a CW space with one cell for every simplex,
including every degenerate simplex: its interior is not identified
with the interior of another cell by a face identification. Each
closed simplex has only finitely many faces, giving closure
finiteness, and the quotient topology is the weak cell topology.
Ordinary realization has one cell for each nondegenerate simplex,
by the preceding unique decomposition. For any simplicial subset,
both induced inclusions are CW-subcomplex inclusions.

Write $K^{(d)}$ for the simplicial subset generated by the
nondegenerate simplices of dimension at most $d$, with
$K^{(-1)}=\varnothing$. Unique decomposition gives the pushout
\begin{equation}\tag{SRC13}
\begin{array}{ccc}
\displaystyle\coprod_{x\in K_d^{\rm nd}}\partial\Delta[d]
 &\longrightarrow& K^{(d-1)}\\
\big\downarrow&&\big\downarrow\\
\displaystyle\coprod_{x\in K_d^{\rm nd}}\Delta[d]
 &\longrightarrow& K^{(d)}
\end{array}
\end{equation}
Both realization constructions preserve this pushout: they are
defined by disjoint unions of products with simplices followed by
the indicated face, or face-and-degeneracy, quotient. In the fat
construction $\|\Delta[d]\|$ is contractible. An $n$-simplex of
$\Delta[d]$ is a nondecreasing list $(v_0,\ldots,v_n)$ with entries
in $\{0,\ldots,d\}$. The explicit contraction is
\begin{equation}\tag{SRC14}
[(v_0,\ldots,v_n);(t_0,\ldots,t_n)]
 \longmapsto
[(v_0,\ldots,v_n,d);
       ((1-u)t_0,\ldots,(1-u)t_n,u)],\qquad 0\le u\le1.
\end{equation}
It respects every face identification, starts at the given point,
and ends at vertex $d$. Ordinary $|\Delta[d]|=\Delta^d$ is also
contractible. Thus $q_{\Delta[d]}$ is a homotopy equivalence.

We recall the elementary gluing argument used with (SRC13). A
pushout of a span $C\leftarrow A\hookrightarrow B$, whose indicated
leg is a CW-subcomplex inclusion, is homotopy equivalent to its
double mapping cylinder
$C\cup(A\times[0,1])\cup B$. To see this, the cylinder collapse
maps to the pushout. The homotopy extension property for $A\subset B$
extends the path on $A$ from the $B$ end of the cylinder to its
$C$ end, starting with the inclusion of $B$. The resulting end
map on $B$ agrees with the inclusion of $C$ along $A$ and defines
a map back from the pushout. The collapse followed by this map
is homotopic to the identity by that extended homotopy on $B$,
the constant homotopy on $C$, and
$(a,s)\mapsto(a,(1-u)s)$ on the connecting cylinder. The reverse
composite is homotopic to the identity after collapsing the same
homotopy. An objectwise homotopy equivalence of two such spans
induces a homotopy equivalence of their double mapping cylinders:
choose inverses on $A,B,C$; their commutation with each attaching
map up to homotopy follows by inserting the chosen inverse
homotopies into the commuting square. On a cylinder use, in order,
the first such attaching homotopy, the cylinder on the inverse
map of $A$, and the reverse of the second attaching homotopy.
This defines an inverse. Its two composites are homotopic to the
identity by the chosen inverse homotopies and cancellation of each
path with its reverse. All pieces agree at the cylinder ends.
This proves the gluing assertion needed here.

Induct on $d$, simultaneously for all simplicial sets of
nondegenerate dimension at most $d$. The assertion for $d=-1$
is empty. In (SRC13) the boundary $\partial\Delta[d]$ and the
previous skeleton have nondegenerate dimension at most $d-1$,
so their comparison maps are homotopy equivalences by induction.
The comparison on each $\Delta[d]$ is a homotopy equivalence by
(SRC14), and the boundary inclusions are CW-subcomplex inclusions.
The gluing argument proves that $q_{K^{(d)}}$ is a homotopy
equivalence. This includes $d=0$, when the boundary is empty.

Finally both full realizations are unions of these simplicial
skeleta. A compact subset of a CW space lies in a finite
subcomplex: if it met infinitely many open cells, selecting one
point in each of infinitely many different cells would produce
an infinite closed discrete subset. Indeed each closed cell meets
only finitely many selected points by closure finiteness; the
same holds for every subset of the selected points, so the weak
topology makes every such subset closed. This contradicts
compactness. Each finite collection of fat cells has a finite
maximum nondegenerate-ancestor dimension. Therefore every sphere,
disk and homotopy in either full realization is contained in
some displayed simplicial skeleton. The finite-stage homotopy
equivalences give isomorphisms on all homotopy groups, including
components, in the union. This proves that (SRC12) is a weak
equivalence without invoking the CW Whitehead theorem. The cited
CW Whitehead theorem then gives the asserted
homotopy equivalence. Naturality of (SRC12) is immediate from its
definition as the additional quotient.

\subsection{The full conclusion and its exact limitations}

For every original divisor $D$ and every finite pointed level $r_+$
there is a commuting square
\begin{equation}\tag{SRC15}
\begin{array}{ccc}
\|X_r\| &\xrightarrow{\|p_r\|}& \|Y_r\|\\
q_{X_r}\big\downarrow&&\big\downarrow q_{Y_r}\\
|X_r| &\xrightarrow{|p_r|}& |Y_r|
\end{array}
\end{equation}
The upper map is a homotopy equivalence by (SRC6)--(SRC11), and
both vertical maps are weak equivalences by (SRC12)--(SRC14).
It follows on every homotopy group and on components that
$|p_r|$ is a weak equivalence; since its source and target are
CW spaces, it is a homotopy equivalence. The maps $|p_r|$ themselves
are pointed and commute with \emph{every} pointed finite-set map.
Thus they give a $\Gamma$-natural levelwise weak equivalence of
the full supported and original restricted objects.

The displayed inverse $\|t_r\|$ and the displayed prism are
levelwise and need not be pointed or $\Gamma$-natural. At an
empty output fibre a pointed finite-set map inserts $\tau$,
whereas $t_r$ would have inserted top support. This is an explicit
failure of naturality, not an omitted check. In fact there is
no pointed simplicial $\Gamma$-section of $p_D$: at degree zero
such a section would give a pointed additive section
$\mathbb A\to G_L(\mathbb A)$. The forced lifts of $x\ne0$
and $-x$ add to $e$, whereas the pointed lift of zero is $\tau$.
The fold map would therefore fail. This does not contradict the
proved $\Gamma$-natural weak equivalence, which concerns the
amplitude projection itself and asserts no natural strict inverse.

In particular the singleton strict spherical set at $\tau^r$ in
ASC8 cannot be used to conclude that $\pi_1(|X_r|,\tau^r)$ is
zero. The exact map identifies that group with
$\pi_1(|Y_r|,0)$, and likewise identifies every higher homotopy
group at the indicated basepoint. As a concrete retained loop,
an admitted signed vector $q\in(N\mathbb Z)^r$ with
$\sum_j|q_j|\le a$ gives the loop $(q,-q)$ at the top-labelled
zero vertex. The edge $c$ from $\tau^r$ to that vertex, followed
by this loop and the reverse geometric path along $c$, is a
based topological loop at $\tau^r$. Under $|p_r|$ it is exactly
the original amplitude loop, since the two $c$ paths become
degenerate. No assertion about that loop's nontriviality is needed
or made without the corresponding calculation in $Y_r$.

All original divisor parameters, finite ideals, rational and
adelic amplitudes, independent support labels and total real
restrictions have been retained in the calculation. Realization
homotopy does not retain every distinction of the strict support
relations: ASC6--ASC13 record those distinctions and their exact
fibres. The result here proves their relation to the original
realization; it does not identify either realization with the
arithmetic spectral cohomology occurring in a Weil trace formula.
